% PDF page 31

\noindent\textit{Proof of Step 3.}
Let $I^n \subset U$ be a cube with edge $\delta$. If $k$ is sufficiently large ($k > n/p - 1$ to be precise) we will prove that $f(C_k \cap I^n)$ has measure zero. Since $C_k$ can be covered by countably many such cubes, this will prove that $f(C_k)$ has measure zero.

From Taylor's theorem, the compactness of $I^n$, and the definition of $C_k$, we see that
\[
f(x + h) = f(x) + R(x, h)
\]
where
\[
\text{1)} \quad \|R(x, h)\| \leq c \|h\|^{k+1}
\]
for $x \in C_k \cap I^n$, $x + h \in I^n$. Here $c$ is a constant which depends only on $f$ and $I^n$. Now subdivide $I^n$ into $r^n$ cubes of edge $\delta/r$. Let $I_1$ be a cube of the subdivision which contains a point $x$ of $C_k$. Then any point of $I_1$ can be written as $x + h$, with
\[
\text{2)} \quad \|h\| \leq \sqrt{n}(\delta/r).
\]
From 1) it follows that $f(I_1)$ lies in a cube of edge $a/r^{k+1}$ centered about $f(x)$, where $a = 2c(\sqrt{n}\,\delta)^{k+1}$ is constant. Hence $f(C_k \cap I^n)$ is contained in a union of at most $r^n$ cubes having total volume
\[
V \leq r^n (a/r^{k+1})^p = a^p r^{n-(k+1)p}.
\]
If $k + 1 > n/p$, then evidently $V$ tends to $0$ as $r \to \infty$; so $f(C_k \cap I^n)$ must have measure zero. This completes the proof of Sard's theorem.


% PDF page 32

\section{The Degree Modulo 2 of a Mapping}

Consider a smooth map $f : S^n \to S^n$. If $y$ is a regular value, recall that $\#f^{-1}(y)$ denotes the number of solutions $x$ to the equation $f(x) = y$. We will prove that the residue class modulo $2$ of $\#f^{-1}(y)$ does not depend on the choice of the regular value $y$. This residue class is called the \emph{mod 2 degree} of $f$.

More generally this same definition works for any smooth map
\[
f : M \to N
\]
where $M$ is compact without boundary, $N$ is connected, and both manifolds have the same dimension. (We may as well assume also that $N$ is compact without boundary, since otherwise the mod 2 degree would necessarily be zero.) For the proof we introduce two new concepts.

\subsection{Smooth Homotopy and Smooth Isotopy}

Given $X \subset \mathbb{R}^k$, let $X \times [0, 1]$ denote the subset\footnote{If $M$ is a smooth manifold without boundary, then $M \times [0, 1]$ is a smooth manifold bounded by two ``copies'' of $M$. Boundary points of $M$ will give rise to ``corner'' points of $M \times [0, 1]$.} of $\mathbb{R}^{k+1}$ consisting of all $(x, t)$ with $x \in X$ and $0 \leq t \leq 1$. Two mappings
\[
f, g: X \to Y
\]
are called \emph{smoothly homotopic} (abbreviated $f \sim g$) if there exists a


% PDF page 33

smooth map $F : X \times [0, 1] \to Y$ with
\[
F(x, 0) = f(x), \qquad F(x, 1) = g(x)
\]
for all $x \in X$. This map $F$ is called a \emph{smooth homotopy} between $f$ and $g$.

Note that the relation of smooth homotopy is an equivalence relation. To see that it is transitive we use the existence of a smooth function $\varphi : [0, 1] \to [0, 1]$ with
\[
\varphi(t) = 0 \quad \text{for} \quad 0 \leq t \leq \tfrac{1}{3}, \qquad
\varphi(t) = 1 \quad \text{for} \quad \tfrac{2}{3} \leq t \leq 1.
\]
(For example, let $\varphi(t) = \lambda(t - \tfrac{1}{3}) / (\lambda(t - \tfrac{1}{3}) + \lambda(\tfrac{2}{3} - t))$, where $\lambda(\tau) = 0$ for $\tau \leq 0$ and $\lambda(\tau) = \exp(-\tau^{-1})$ for $\tau > 0$.) Given a smooth homotopy $F$ between $f$ and $g$, the formula $G(x, t) = F(x, \varphi(t))$ defines a smooth homotopy $G$ with
\[
G(x, t) = f(x) \quad \text{for} \quad 0 \leq t \leq \tfrac{1}{3}, \qquad
G(x, t) = g(x) \quad \text{for} \quad \tfrac{2}{3} \leq t \leq 1.
\]
Now if $f \sim g$ and $g \sim h$, then, with the aid of this construction, it is easy to prove that $f \sim h$.

If $f$ and $g$ happen to be diffeomorphisms from $X$ to $Y$, we can also define the concept of a ``smooth isotopy'' between $f$ and $g$. This also will be an equivalence relation.

\textsc{Definition.} The diffeomorphism $f$ is \emph{smoothly isotopic} to $g$ if there exists a smooth homotopy $F : X \times [0, 1] \to Y$ from $f$ to $g$ so that, for each $t \in [0, 1]$, the correspondence $x \mapsto F(x, t)$ maps $X$ diffeomorphically onto $Y$.

It will turn out that the mod 2 degree of a map depends only on its smooth homotopy class:

\noindent\textbf{Homotopy Lemma.} \textit{Let $f, g : M \to N$ be smoothly homotopic maps between manifolds of the same dimension, where $M$ is compact and without boundary. If $y \in N$ is a regular value for both $f$ and $g$, then}
\[
\#f^{-1}(y) \equiv \#g^{-1}(y) \pmod{2}.
\]

\noindent\textit{Proof.}
Let $F : M \times [0, 1] \to N$ be a smooth homotopy between $f$ and $g$. First suppose that $y$ is also a regular value for $F$. Then $F^{-1}(y)$

